
\subsection{Ortogonalių matricų grupė $\Orth3$}

\begin{definition}
	Ortogonalių matricų grupė $\Orth3$ yra sudaryta iš $3\times3$ matricų
	$R$ virš realiųjų skaičių, tokių kad, lygybė 
	\begin{equation}
		(R \vb v, R \vb w) = (\vb v, \vb w)
	\end{equation}
	yra tenkinama su visais vektoriais $\vb v$ ir $\vb w$.
\end{definition}

Ortogonaliomis matricomis jos yra vadinamos, nes bet kokie ortogonalūs
vektoriai $\vb v$ ir $\vb w$ yra atvaizduojami į ortogonalius vektorius 
$R\vb v$ ir $R\vb w$. Šios matricos taip pat pasižymi tuo, kad matricą $R$ sudauginus
su savo transponuotą $R^t$ rezultatas yra vienetinę matricą $I$, tai
\begin{equation}
	\begin{split}
		\det(RR^t) = \det(I)
		\\
		\qty(\det(R))^2 = 1
		\\
		\det(R) = \pm 1
	\end{split}
\end{equation}
Determinantas gerbia matricų daugybą todėl tai yra grupės homomorfizmas
iš $\Orth3$ į $\qty{1,-1}$. To pasekoje $\Orth3$ Yra padalijama į dvi dalis:
teisingus sukimus (\textit{ang.} proper rotations), kuriems $\det(R) = 1$, ir
neteisingus sukimus (\textit{ang.} improper rotations), kuriems $\det(R) = -1$.
\begin{example}
	\leavevmode
	\begin{itemize}
		\item 
			$\smqty(
			\cos(\alpha) & -\sin(\alpha) &   \\
			\sin(\alpha) &  \cos(\alpha) &   \\
				     &		     & 1
			)$
			yra teisingo sukimo matrica, kuri pasuka vektorių
			kampu $\alpha$ aplink $z$ ašį;
		\item
			$\smqty(
			1 &   &   \\
			  & 1 &   \\
			  &   & -1
			)$ yra neteisingo sukimo matrica, kuri vektorių atspindi
			per $xy$ plokštumą;
		\item
			$\smqty(
			-1 &    &   \\
			   & -1 &   \\
			   &    & -1
			)$ yra neteisingo sukimo matrica, kuri vadinama
			inversijos transformacija ir ji vektorių siunčia į savo
			minusą.
	\end{itemize}
\end{example}
Teisingus sukimus galime realizuoti kaip tolydų objekto pasukimą. Neteisingi
pasukimai visata turės savyje atspindžio arba inversijos veiksmą, todėl jų
neįmanoma pasiekti tolydžiu pasukimu. Pasukimai be atspindžio ar inversijos
sudaro $\Orth3$ pogrupį vadinamą $\SO3$.
\begin{definition}
	Specialioji ortogonalių matricų grupė žymima $\SO3$ ir yra sudaryta iš
	$3\times3$ matricų $R$ virš realiųjų skaičių, tokių kad,
	\begin{itemize}
		\item lygybė $(R \vb v, R \vb w) = (\vb v, \vb w)$ yra
			tenkinama su visais vektoriais $\vb v$ ir $\vb w$
		\item ir $\det(R) = 1$.
	\end{itemize}
\end{definition}

\begin{preposition}
	$\Orth3$ yra izomorfiška grupių sandaugai $\SO3 \times \qty{1,-1}$.
\end{preposition}
\begin{proof}
	Tegul $R' \in \Orth3$ ir $\det(R')=-1$, tada $R'=IR'=PPR'$,
	čia $P$ yra inversijos matrica $\smqty(
			-1 &    &   \\
			   & -1 &   \\
			   &    & -1
	)$ ir $PR' \in \SO3$, nes 
	\begin{equation}
		\det(PR')=\det(P)\det(R')=(-1)(-1)=1 \,.
	\end{equation}
	Kiekvieną $\Orth3$ elementą, kuris nepriklauso $\SO3$ galima išskaidyti
	į $PR$, kur $R \in \SO3$. Tada izomorfizmas yra duodamas funkcijos
	\begin{equation}
		\begin{split}
			f: \SO3 \times \qty{1,-1} &\to \Orth3\,,\\
			(R,1) &\mapsto R \\
			(R,-1) &\mapsto PR \,.
		\end{split}
	\end{equation}
	Unikalus išskaidymas garantuoja, kad $f^{-1}$ egzistuoja. Kad įrodyti
	kad $f$ tikrai yra izomorfizmas užtenka parodyti, kad $f$ išlaiko grupės
	struktūrą dauginant teisingus sukimus, teisingą sukimą su neteisingu ir
	dauginant du neteisingus sukimus.
	\begin{equation}
		\begin{split}
			f((R_1,1))f((R_2,1))
			&= R_1 R_2 \\
			&= f((R_1 R_2,1)) \\
			&= f((R_1,1)(R_2,1)) \,,
		\end{split}
	\end{equation}
	\begin{equation}
		\begin{split}
			f((R_1,1))f((R_2,-1))
			&= R_1 P R_2 \\
			&= P R_1 R_2 \\
			&= f((R_1 R_2,-1)) \\
			&= f((R_1,1)(R_2,-1)) \,,
		\end{split}
	\end{equation}
	\begin{equation}
		\begin{split}
			f((R_1,-1))f((R_2,-1))
			&= P R_1 P R_2 \\
			&= P^2 R_1 R_2 \\
			&= R_1 R_2 \\
			&= f((R_1 R_2,1)) \\
			&= f((R_1,-1)(R_2,-1)) \,.
		\end{split}
	\end{equation}
	Čia pasinaudojome faktu kad $PR=RP$ su visais $R \in \SO3$ ir $P^2=I$. 
	Egzistuoja izomorfizmas $f$, tad $\Orth3$ ir $\SO3 \times \qty{1,-1}$
	yra izomorfiškos grupės.
\end{proof}




